\chapter{SSH, 鍵認証, ファイル転送}
\label{app:ssh}
\index{ssh@SSH|textbf}
\index{かぎにんしょう@鍵認証|textbf}

\section{サーバーへの SSH 接続と鍵認証}

SSH（Secure Shell）は、離れた計算機との通信を暗号化し、その計算機上でコマンドを実行するための仕組みである。研究室や計算サーバーに接続し、大規模なデータ処理や長時間ジョブを実行する場合に利用する。GitHub への SSH 接続も同じ認証方式を用いるが、GitHub は対話的なシェルの接続先ではなく、Git 操作の認証先である。

\subsection{サーバーへの接続}

初回の SSH 接続では、接続先のホスト鍵の指紋（fingerprint）を確認することがある。これは接続先サーバーを識別するためのもので、自分の認証用公開鍵とは異なる。管理者など信頼できる経路から得た指紋と照合して登録する。ホスト鍵が変わったという警告では、再構築などの理由を確認してから更新する。詳細は OpenSSH の資料~\cite{openssh_manual} を参照する。
\index{ssh command@\code{ssh}|textbf}

最も基本的な形は次である。

\begin{lstlisting}[numbers=none, xleftmargin=2\zw]
$ ssh user@server.example.jp
\end{lstlisting}

\code{user} はサーバー上の自分のユーザー名、\code{server.example.jp} はホスト名である。ポート番号が既定の 22 でないなら \code{-p} を付ける。

\begin{lstlisting}[numbers=none, xleftmargin=2\zw]
$ ssh -p 2222 user@server.example.jp
\end{lstlisting}

接続後に入力したコマンドは、原則として接続先の計算機上で実行される。表示されるファイルも接続先にあるため、ローカルファイルとの取り違えは意図しない変更や削除につながる。作業を始める前に \code{hostname}、\code{whoami}、\code{pwd} を実行すれば、接続先、利用者名、作業ディレクトリを確認できる。

\subsection{SSH 鍵の作成}
\index{ssh-keygen@\code{ssh-keygen}|textbf}
\index{こうかいかぎ@公開鍵|textbf}
\index{ひみつかぎ@秘密鍵|textbf}

公開鍵認証では、公開鍵と秘密鍵の対を使用する。公開鍵をサーバーへ登録し、秘密鍵は接続元の計算機から外部へ渡さない。

\begin{lstlisting}[numbers=none, xleftmargin=2\zw]
$ ssh-keygen -t ed25519 -C "your_email@example.com"
\end{lstlisting}

既定の保存先を選ぶと、\code{\textasciitilde/.ssh/id\_ed25519} が秘密鍵、\code{\textasciitilde/.ssh/id\_ed25519.pub} が公開鍵になる。サーバーへ登録するのは \code{.pub} が付いた公開鍵だけである。\code{.pub} が付かない秘密鍵は共有しない。

\begin{lstlisting}[numbers=none, xleftmargin=2\zw]
$ ls -l ~/.ssh
\end{lstlisting}

秘密鍵の権限が緩いと SSH が拒否することがある。その場合は次のように直す。

\begin{lstlisting}[numbers=none, xleftmargin=2\zw]
$ chmod 700 ~/.ssh
$ chmod 600 ~/.ssh/id_ed25519
$ chmod 644 ~/.ssh/id_ed25519.pub
\end{lstlisting}

\subsection{サーバーへの公開鍵登録}
\index{ssh-copy-id@\code{ssh-copy-id}|textbf}

\code{ssh-copy-id} は、公開鍵を接続先の \code{authorized\_keys} へ追加するコマンドである。

\begin{lstlisting}[numbers=none, xleftmargin=2\zw]
$ ssh-copy-id user@server.example.jp
\end{lstlisting}

特定の公開鍵を使いたいなら \code{-i} で指定できる。

\begin{lstlisting}[numbers=none, xleftmargin=2\zw]
$ ssh-copy-id -i ~/.ssh/id_ed25519.pub user@server.example.jp
\end{lstlisting}

\code{ssh-copy-id} は、指定した公開鍵をサーバー側の \code{\textasciitilde/.ssh/authorized\_keys} へ追加する。公開鍵の転送と追記を一つのコマンドで実行できる。

\subsubsection{\code{ssh-copy-id} が使えない場合}

\code{ssh-copy-id} が使えない環境では、公開鍵を手動で登録する。登録対象は、次のコマンドで表示される 1 行全体である。

\begin{lstlisting}[numbers=none, xleftmargin=2\zw]
$ cat ~/.ssh/id_ed25519.pub
\end{lstlisting}

表示された 1 行全体を、サーバー側の \code{\textasciitilde/.ssh/authorized\_keys} へ追記する。サーバーへ通常のパスワードログインができるなら、自分でログインして追記してもよい。

\subsection{\code{ssh-agent} による秘密鍵の管理}
\index{ssh-agent@\code{ssh-agent}|textbf}

パスフレーズ付きの鍵を使う場合は、\code{ssh-agent} に読み込ませておくと毎回の入力を減らせる。

\begin{lstlisting}[numbers=none, xleftmargin=2\zw]
$ eval "$(ssh-agent -s)"
$ ssh-add ~/.ssh/id_ed25519
\end{lstlisting}

\section{接続設定と GitHub 連携}

\subsection{\code{\textasciitilde/.ssh/config} による接続設定}
\index{ssh config@\code{\textasciitilde/.ssh/config}|textbf}

\code{\textasciitilde/.ssh/config} には、接続先ごとのホスト名、利用者名、ポート番号、秘密鍵ファイルを記録できる。次の設定では、これらの接続条件に \code{labserver} という名前を付ける。

\begin{lstlisting}[numbers=none, xleftmargin=2\zw, title={\textasciitilde/.ssh/config}]
Host labserver
    HostName server.example.jp
    User student
    Port 2222
    IdentityFile ~/.ssh/id_ed25519
\end{lstlisting}

この設定を保存すると、次のコマンドが同じ接続条件を参照する。

\begin{lstlisting}[numbers=none, xleftmargin=2\zw]
$ ssh labserver
\end{lstlisting}

この設定ファイルには接続先、利用者名、秘密鍵ファイルの場所が記録されるが、秘密鍵そのものは含まれない。ただし、第三者による書き換えを防ぐため、所有者だけが読み書きできる権限を設定する。

\begin{lstlisting}[numbers=none, xleftmargin=2\zw]
$ chmod 600 ~/.ssh/config
\end{lstlisting}

\subsection{GitHub への SSH 鍵登録}

通常のサーバーとは違って、GitHub には \code{ssh-copy-id} を使わない。GitHub はシェル接続先ではなく、アカウント設定に公開鍵を登録する相手だからである。

GitHub の設定画面にある \code{SSH and GPG keys} へ公開鍵を登録する。GitHub CLI \code{gh} を使う場合は、認証手順の中で既存鍵の登録または新規鍵の作成を選択できる。

\begin{lstlisting}[numbers=none, xleftmargin=2\zw]
$ gh auth login --git-protocol ssh --web
\end{lstlisting}

GitHub CLI を使う場合は、既存の鍵を検出してアップロードするか、新しい鍵を作るかを対話的に選べる。GitHub CLI を使わない場合は、\code{cat ~/.ssh/id\_ed25519.pub} で表示した公開鍵を Web 画面へ貼り付ければよい。

\section{\code{scp}・\code{rsync} によるファイル転送}

\code{scp} は指定したファイルまたはディレクトリを転送する。\code{rsync} は送信元と転送先の差分を調べ、除外条件や中断後の部分ファイルの保持を指定できる。単発の転送には \code{scp}、同じディレクトリの反復同期には \code{rsync} を用いる。

\subsection{\code{scp} の最小例}
\index{scp@\code{scp}|textbf}

\code{scp} では、ローカルとリモートのどちらを送信元に書くかによって転送方向が決まる。両方向の最小例を次に示す。

\begin{lstlisting}[numbers=none, xleftmargin=2\zw]
$ scp result.csv user@server.example.jp:~/work/
$ scp user@server.example.jp:~/work/result.csv .
\end{lstlisting}

1 行目はローカルからサーバーへ送信し、2 行目はサーバーからローカルへ受信する。ディレクトリ全体を転送する場合は \code{-r} を付ける。

\subsection{差分同期に \code{rsync} を用いる理由}
\index{rsync@\code{rsync}|textbf}

\code{rsync} は、送信元と転送先の差分を転送し、\code{--partial} で中断時の部分ファイルを保持し、\code{--exclude} で転送対象を除外できる。解析用ディレクトリやコード一式を反復同期する場合は、変更のないファイルも再転送する \code{scp} より転送量を抑えられる。

\begin{lstlisting}[numbers=none, xleftmargin=2\zw]
$ rsync -aP ./analysis/ user@server.example.jp:~/work/analysis/
\end{lstlisting}

主なオプションの意味を次に示す。

\begin{description}
\item[\code{-a}] アーカイブ（archive）モードである。再帰コピーを行い、時刻や権限などもまとめて保つ。
\item[\code{-P}] \code{--partial --progress} の省略形である。途中ファイルを残しつつ進捗を表示する。
\item[\code{-v}] 何を転送しているか詳しく表示する。
\item[\code{-z}] 転送時に圧縮する。低速回線では有利だが、既に圧縮済みの大きなデータでは効果が薄いこともある。
\item[\code{--exclude}] 特定のファイルやディレクトリを転送対象から除外する。
\end{description}

次の例では、アーカイブモードと進捗表示を有効にし、転送先で再作成できる仮想環境、キャッシュ、ビルド成果物を除外する。

\begin{lstlisting}[numbers=none, xleftmargin=2\zw]
$ rsync -aPv \
  --exclude='.venv/' \
  --exclude='__pycache__/' \
  --exclude='build/' \
  --exclude='.mypy_cache/' \
  ./analysis/ user@server.example.jp:~/work/analysis/
\end{lstlisting}

この例では、依存関係の定義やソースコードから転送先で再生成できる仮想環境、キャッシュ、ビルド成果物を除外し、転送量を減らしている。

\subsubsection{末尾のスラッシュによる転送結果の違い}
\index{rsync@\code{rsync}!まつびのすらっしゅ@末尾のスラッシュ}

\code{rsync} では、送信元パスの末尾に \code{/} を付けるかどうかで意味が変わる。

\begin{lstlisting}[numbers=none, xleftmargin=2\zw]
$ rsync -aP analysis user@server:~/work/
$ rsync -aP analysis/ user@server:~/work/
\end{lstlisting}

1 行目は \code{analysis} というディレクトリごと送り、2 行目は \code{analysis} の中身を送る。末尾のスラッシュを取り違えると転送先の階層が 1 段ずれる。初回の転送前には、\code{--dry-run} を加えた実行で転送対象を確認できる。

転送条件やオプションの詳細は rsync の公式マニュアル~\cite{rsync_manual} で確認できる。
